(1) 体自由电荷满足 \(\partial_t\rho_f+\nabla\cdot\boldsymbol J=0\)。用欧姆定律 \(\boldsymbol J=\sigma\boldsymbol E\) 及 \(\nabla\cdot\boldsymbol D=\rho_f\)、\(\boldsymbol D=\varepsilon\boldsymbol E\)，得到 \(\partial_t\rho_f+(\sigma/\varepsilon)\rho_f=0\)。该方程不要求电场在空间均匀，仅要求所讨论介质区域内参数 \(\sigma,\varepsilon\) 为常数。
(2) 每个介质内部点均按同样的一阶常微分方程衰减，故 \(\rho_f(\boldsymbol r,t)=\rho_0(\boldsymbol r)e^{-t/\tau_M}\)，\(\tau_M=\varepsilon/\sigma\)。弛豫时间具有秒的量纲，良导体取极短时间而高绝缘体取较长时间。若同时存在外加驱动电荷注入，方程右端会有源项，不能继续直接采用纯指数无源解。
(3) 对有限体积积分连续性方程得 \(dQ_{\rm bulk}/dt=-\oint_{\partial V}\boldsymbol J\cdot d\boldsymbol S\)，体电荷的减少意味着电荷流向界面或外电极，而不是电荷凭空湮灭。若边界不允许自由电荷穿出系统，初始净自由电荷会转移到表面，满足界面法向电位移跃变 \(\hat n\cdot(\boldsymbol D_{\rm out}-\boldsymbol D_{\rm in})=\sigma_f^{\rm surf}\)。
(4) 含位移电流的安培方程是 \(\nabla\times\boldsymbol H=\boldsymbol J+\partial_t\boldsymbol D\)；两侧取散度得到 \(0=\nabla\cdot\boldsymbol J+\partial_t\rho_f\)，正是连续性方程。若 \(\sigma\) 依赖温度/场强、\(\varepsilon\) 随时间或频率变化，或者存在扩散电流 \(\boldsymbol J=-D\nabla\rho_f+\sigma\boldsymbol E\)，局部动力学不再是上述简单常系数指数律。
